
\documentclass{article}
\usepackage{colortbl}
\usepackage{makecell}
\usepackage{multirow}
\usepackage{supertabular}

\begin{document}

\newcounter{utterance}

\centering \large Interaction Transcript for game `clean\_up', experiment `2\_hard\_5obj\_lt', episode 0 with Qwen3.6{-}35B{-}A3B{-}FP8{-}without{-}reasoning.
\vspace{24pt}

{ \footnotesize  \setcounter{utterance}{1}
\setlength{\tabcolsep}{0pt}
\begin{supertabular}{c@{$\;$}|p{.15\linewidth}@{}p{.15\linewidth}p{.15\linewidth}p{.15\linewidth}p{.15\linewidth}p{.15\linewidth}}
    \# & $\;$A & \multicolumn{4}{c}{Game Master} & $\;\:$B\\
    \hline

    \theutterance \stepcounter{utterance}  
    & & \multicolumn{4}{p{0.6\linewidth}}{
        \cellcolor[rgb]{0.9,0.9,0.9}{
            \makecell[{{p{\linewidth}}}]{
                \texttt{\tiny{[P1$\langle$GM]}}
                \texttt{Aš esu jūsų žaidimo vedėjas, o jūs žaidžiate bendrą žaidimą, kurio žaidimo lenta yra tokia lentelė:} \\
\\ 
\texttt{```} \\
\texttt{ 1234567} \\
\texttt{╔══╤════╗} \\
\texttt{║◌◌│◌◌◌H║ 1} \\
\texttt{╟──┤◌◌┌─╢ 2} \\
\texttt{║CI├──┤◌║ 3} \\
\texttt{╟──┘◌◌├─╢ 4} \\
\texttt{║◌◌◌◌◌│◌║ 5} \\
\texttt{╟─┐W┌─┼─╢ 6} \\
\texttt{║◌│T│◌│◌║ 7} \\
\texttt{╚═╧═╧═╧═╝} \\
\\ 
\texttt{```} \\
\\ 
\texttt{* Viršutiniame krašte x koordinatės didėja į dešinę, o dešiniajame krašte y koordinatės didėja žemyn.} \\
\texttt{* Tinklelyje atsitiktine tvarka išdėstyti šie objektai: 'W', 'I', 'T', 'C', 'H'.} \\
\\ 
\\ 
\texttt{Kitas žaidėjas mato kitą žaidimo lentos variantą, kuriame objektai išdėstyti skirtingose atsitiktinėse vietose. Jūs negalite matyti kito žaidėjo lentos, o jis negali matyti jūsų lentos.} \\
\\ 
\\ 
\texttt{**Tikslas:**} \\
\\ 
\texttt{Abu žaidėjai turi perkelti objektus savo fone taip, kad vienodi objektai atsidurtų tose pačiose koordinatėse. Turite bendrauti su kitu žaidėju ir susitarti dėl bendros tikslinės konfigūracijos.} \\
\\ 
\\ 
\texttt{**Taisyklės:**} \\
\\ 
\texttt{* Kiekvieno ėjimo metu galite išsiųsti lygiai vieną iš šių dviejų komandų:} \\
\texttt{ 1. `SAKYTI: <ŽINUTĖ>`: norėdami išsiųsti žinutę (viską iki kitos eilutės pertraukos) kitam žaidėjui. Aš ją persiųsiu jūsų partneriui.} \\
\texttt{2. `JUDĖTI: <OBJEKTAS>, (<X>, <Y>)`: perkelti objektą į naują poziciją, kur `<X>` yra stulpelis, o `<Y>` yra eilutė. Aš jums pranešiu, ar jūsų perkėlimas buvo sėkmingas, ar ne.} \\
\texttt{* Jei nesilaikysite formato arba siųsite kelias komandas iš karto, turėsiu jus nubausti.} \\
\texttt{* Jei abu žaidėjai sukaups daugiau nei 12 baudų, abu pralaimi žaidimą.} \\
\texttt{* Labai svarbu, kad bendrautumėte su kitu žaidėju dėl savo tikslinės būsenos! Vienintelis būdas, kuriuo galite perduoti savo strategiją kitam žaidėjui, yra komanda `SAKYTI: <ŽINUTĖ>`!} \\
\\ 
\\ 
\texttt{**Objektų perkėlimas**} \\
\\ 
\texttt{* Objektus galima perkelti tik į langelius, esančius tinklelio ribose. Tikslinis langelis turi būti tuščias, t. y. jame turi būti tik simbolis '◌'.} \\
\texttt{* Jei bandysite perkelti objektą į vietą, kuri nėra tuščia, arba bandysite jį perkelti už tinklelio ribų, turėsiu jus nubausti. Gausite dar vieną bandymą.} \\
\texttt{* Prieš atlikdami perkėlimą, dukart patikrinkite, ar tikslinė vieta yra tuščia ir ar joje nėra jokios raidės, rėmelio ar linijos!} \\
\\ 
\\ 
\texttt{**Žaidimo pabaiga**} \\
\\ 
\texttt{Jei manote, kad pasiekėte tikslą {-} sulygiavote visus objektus, galite paprašyti kito žaidėjo baigti žaidimą siųsdami `SAKYTI: baigta?`. Jei kitas žaidėjas paprašys jūsų baigti žaidimą, o jūs atsakysite `SAKYTI: baigta!`, žaidimas baigsis.} \\
\\ 
\texttt{Abu žaidėjai laimi, jei žaidimas baigiasi per 20 raundus, kur vienas raundas apibrėžiamas kaip du žaidėjai, kiekvienas išsiuntęs po vieną galiojančią komandą.} \\
\\ 
\texttt{**Taškų skaičiavimas:**} \\
\\ 
\texttt{Kuo arčiau vienodų objektų abiejose žaidimo lentose, tuo daugiau taškų gausite. Baudos mažina jūsų taškus. Ar pavyks pagerinti rekordą?} \\
\\ 
\\ 
\texttt{Išsiųskite žinutę kitam žaidėjui, kad pradėtumėte žaidimą!} \\
            }
        }
    }
    & \\ \\

    \theutterance \stepcounter{utterance}  
    & \multicolumn{4}{p{0.6\linewidth}}{
        \cellcolor[rgb]{0.9,0.9,0.9}{
            \makecell[{{p{\linewidth}}}]{
                \texttt{\tiny{[P1$\rangle$GM]}}
                \texttt{SAKYTI: Sveiki! Kadangi mes negalime matyti vienas kito lentų, mums reikia susitarti dėl bendros tikslinės konfigūracijos. Kadangi objektai pas mus išdėstyti skirtingai, aš pasiūlyčiau vieną iš dviejų strategijų:} \\
\\ 
\texttt{1. **Standartizuota tvarka:** Mes abu sutariame, kad objektai bus išdėstyti tam tikra tvarka (pvz., pagal raidžių aibę: C, H, I, T, W) užimant konkrečias tuščias vietas.} \\
\texttt{2. **Mano siūlymas:** Aš matau savo lentą. Pabandykime susitarti dėl to, kad mes abu perkelsime objektus į **kairę viršutinę dalį**, užimdami pirmas laisvas vietas eilutėmis.} \\
\\ 
\texttt{Kadangi aš negaliu žinoti, kur jūsų objektai yra, man reikia, kad **jūs** man pasakytumėte, kokias koordinates (X, Y) jūsų lentos tuščiuose langeliuose (kur yra '◌') jūs norėtumėte naudoti kaip tikslines vietas mūsų objektams.} \\
\\ 
\texttt{Prašau, parašykite man sąrašą iš 5 koordinatų (X, Y), kuriose jūsų lentos yra laisvi langeliai ('◌'). Aš tuos pačius koordinates naudosiu ir savo lentos tuščiuose langeliuose. Tada mes abu perkelsime savo objektus į tas pačias koordinates, ir jie sutaps.} \\
\\ 
\texttt{Pavyzdžiui, jei manote, kad geriausia naudoti šias vietas: (1,1), (2,1), (3,1), (4,1), (5,1) – parašykite tai. Arba pasirinkite kitas laisvas vietas, kurios yra patogios abiem.} \\
\\ 
\texttt{Laukiu jūsų pasiūlymo dėl koordinatžių!} \\
            }
        }
    }
    & & \\ \\

    \theutterance \stepcounter{utterance}  
    & & & \multicolumn{2}{p{0.3\linewidth}}{
        \cellcolor[rgb]{0.9,0.9,0.9}{
            \makecell[{{p{\linewidth}}}]{
                \texttt{\tiny{[GM$|$GM]}}
                \texttt{SAKYTI: Sveiki! Kadangi mes negalime matyti vienas kito lentų, mums reikia susitarti dėl bendros tikslinės konfigūracijos. Kadangi objektai pas mus išdėstyti skirtingai, aš pasiūlyčiau vieną iš dviejų strategijų:} \\
\\ 
\texttt{1. **Standartizuota tvarka:** Mes abu sutariame, kad objektai bus išdėstyti tam tikra tvarka (pvz., pagal raidžių aibę: C, H, I, T, W) užimant konkrečias tuščias vietas.} \\
\texttt{2. **Mano siūlymas:** Aš matau savo lentą. Pabandykime susitarti dėl to, kad mes abu perkelsime objektus į **kairę viršutinę dalį**, užimdami pirmas laisvas vietas eilutėmis.} \\
\\ 
\texttt{Kadangi aš negaliu žinoti, kur jūsų objektai yra, man reikia, kad **jūs** man pasakytumėte, kokias koordinates (X, Y) jūsų lentos tuščiuose langeliuose (kur yra '◌') jūs norėtumėte naudoti kaip tikslines vietas mūsų objektams.} \\
\\ 
\texttt{Prašau, parašykite man sąrašą iš 5 koordinatų (X, Y), kuriose jūsų lentos yra laisvi langeliai ('◌'). Aš tuos pačius koordinates naudosiu ir savo lentos tuščiuose langeliuose. Tada mes abu perkelsime savo objektus į tas pačias koordinates, ir jie sutaps.} \\
\\ 
\texttt{Pavyzdžiui, jei manote, kad geriausia naudoti šias vietas: (1,1), (2,1), (3,1), (4,1), (5,1) – parašykite tai. Arba pasirinkite kitas laisvas vietas, kurios yra patogios abiem.} \\
\\ 
\texttt{Laukiu jūsų pasiūlymo dėl koordinatžių!} \\
            }
        }
    }
    & & \\ \\

    \theutterance \stepcounter{utterance}  
    & & & \multicolumn{2}{p{0.3\linewidth}}{
        \cellcolor[rgb]{0.9,0.9,0.9}{
            \makecell[{{p{\linewidth}}}]{
                \texttt{\tiny{[GM$|$GM]}}
                \texttt{Invalid format: tail is not empty: '} \\
\\ 
\texttt{1. **Standartizuota tvarka:** Mes abu sutariame, kad objektai bus išdėstyti tam tikra tvarka (pvz., pagal raidžių aibę: C, H, I, T, W) užimant konkrečias tuščias vietas.} \\
\texttt{2. **Mano siūlymas:** Aš matau savo lentą. Pabandykime susitarti dėl to, kad mes abu perkelsime objektus į **kairę viršutinę dalį**, užimdami pirmas laisvas vietas eilutėmis.} \\
\\ 
\texttt{Kadangi aš negaliu žinoti, kur jūsų objektai yra, man reikia, kad **jūs** man pasakytumėte, kokias koordinates (X, Y) jūsų lentos tuščiuose langeliuose (kur yra '◌') jūs norėtumėte naudoti kaip tikslines vietas mūsų objektams.} \\
\\ 
\texttt{Prašau, parašykite man sąrašą iš 5 koordinatų (X, Y), kuriose jūsų lentos yra laisvi langeliai ('◌'). Aš tuos pačius koordinates naudosiu ir savo lentos tuščiuose langeliuose. Tada mes abu perkelsime savo objektus į tas pačias koordinates, ir jie sutaps.} \\
\\ 
\texttt{Pavyzdžiui, jei manote, kad geriausia naudoti šias vietas: (1,1), (2,1), (3,1), (4,1), (5,1) – parašykite tai. Arba pasirinkite kitas laisvas vietas, kurios yra patogios abiem.} \\
\\ 
\texttt{Laukiu jūsų pasiūlymo dėl koordinatžių!} \\
            }
        }
    }
    & & \\ \\

    \theutterance \stepcounter{utterance}  
    & & \multicolumn{4}{p{0.6\linewidth}}{
        \cellcolor[rgb]{0.9,0.9,0.9}{
            \makecell[{{p{\linewidth}}}]{
                \texttt{\tiny{[P1$\langle$GM]}}
                \texttt{Nuobauda: Jūsų žinutėje neturi būti nieko po komandos!} \\
\\ 
\texttt{Įsitikinkite, kad jūsų atsakyme yra tik `SAKYTI: <ŽINUTĖ>` arba `JUDĖTI: <OBJEKTAS>, (<X>, <Y>)`, ir nieko daugiau!} \\
\\ 
\texttt{Jūs bendrai sukaupėte 1 iš 12 baudų.} \\
\texttt{Pabandykite dar kartą!} \\
            }
        }
    }
    & \\ \\

    \theutterance \stepcounter{utterance}  
    & \multicolumn{4}{p{0.6\linewidth}}{
        \cellcolor[rgb]{0.9,0.9,0.9}{
            \makecell[{{p{\linewidth}}}]{
                \texttt{\tiny{[P1$\rangle$GM]}}
                \texttt{SAKYTI: Sveiki! Kadangi mūsų lentos skiriasi, mums reikia susitarti dėl bendrų tikslinių pozicijų. Prašau, pasiūlykite 5 konkrečias (X, Y) koordinates, kuriose jūsų lentos yra laisvi langeliai ('◌'). Aš tuos pačius koordinates naudosiu savo lentos laisvoms vietoms perkelti savo objektus. Tada mūsų objektai sutaps. Laukiu jūsų koordinatžių!} \\
            }
        }
    }
    & & \\ \\

    \theutterance \stepcounter{utterance}  
    & & & \multicolumn{2}{p{0.3\linewidth}}{
        \cellcolor[rgb]{0.9,0.9,0.9}{
            \makecell[{{p{\linewidth}}}]{
                \texttt{\tiny{[GM$|$GM]}}
                \texttt{SAKYTI: Sveiki! Kadangi mūsų lentos skiriasi, mums reikia susitarti dėl bendrų tikslinių pozicijų. Prašau, pasiūlykite 5 konkrečias (X, Y) koordinates, kuriose jūsų lentos yra laisvi langeliai ('◌'). Aš tuos pačius koordinates naudosiu savo lentos laisvoms vietoms perkelti savo objektus. Tada mūsų objektai sutaps. Laukiu jūsų koordinatžių!} \\
            }
        }
    }
    & & \\ \\

    \theutterance \stepcounter{utterance}  
    & & \multicolumn{4}{p{0.6\linewidth}}{
        \cellcolor[rgb]{0.9,0.9,0.9}{
            \makecell[{{p{\linewidth}}}]{
                \texttt{\tiny{[GM$\rangle$P2]}}
                \texttt{Aš esu jūsų žaidimo vedėjas, o jūs žaidžiate bendrą žaidimą, kurio žaidimo lenta yra tokia lentelė:} \\
\\ 
\texttt{```} \\
\texttt{ 1234567} \\
\texttt{╔══╤════╗} \\
\texttt{║◌W│◌◌◌◌║ 1} \\
\texttt{╟──┤◌H┌─╢ 2} \\
\texttt{║◌◌├──┤◌║ 3} \\
\texttt{╟──┘◌◌├─╢ 4} \\
\texttt{║◌◌◌◌◌│◌║ 5} \\
\texttt{╟─┐C┌─┼─╢ 6} \\
\texttt{║I│◌│T│◌║ 7} \\
\texttt{╚═╧═╧═╧═╝} \\
\\ 
\texttt{```} \\
\\ 
\texttt{* Viršutiniame krašte x koordinatės didėja į dešinę, o dešiniajame krašte y koordinatės didėja žemyn.} \\
\texttt{* Tinklelyje atsitiktine tvarka išdėstyti šie objektai: 'W', 'I', 'T', 'C', 'H'.} \\
\\ 
\\ 
\texttt{Kitas žaidėjas mato kitą žaidimo lentos variantą, kuriame objektai išdėstyti skirtingose atsitiktinėse vietose. Jūs negalite matyti kito žaidėjo lentos, o jis negali matyti jūsų lentos.} \\
\\ 
\\ 
\texttt{**Tikslas:**} \\
\\ 
\texttt{Abu žaidėjai turi perkelti objektus savo fone taip, kad vienodi objektai atsidurtų tose pačiose koordinatėse. Turite bendrauti su kitu žaidėju ir susitarti dėl bendros tikslinės konfigūracijos.} \\
\\ 
\\ 
\texttt{**Taisyklės:**} \\
\\ 
\texttt{* Kiekvieno ėjimo metu galite išsiųsti lygiai vieną iš šių dviejų komandų:} \\
\texttt{ 1. `SAKYTI: <ŽINUTĖ>`: norėdami išsiųsti žinutę (viską iki kitos eilutės pertraukos) kitam žaidėjui. Aš ją persiųsiu jūsų partneriui.} \\
\texttt{2. `JUDĖTI: <OBJEKTAS>, (<X>, <Y>)`: perkelti objektą į naują poziciją, kur `<X>` yra stulpelis, o `<Y>` yra eilutė. Aš jums pranešiu, ar jūsų perkėlimas buvo sėkmingas, ar ne.} \\
\texttt{* Jei nesilaikysite formato arba siųsite kelias komandas iš karto, turėsiu jus nubausti.} \\
\texttt{* Jei abu žaidėjai sukaups daugiau nei 12 baudų, abu pralaimi žaidimą.} \\
\texttt{* Labai svarbu, kad bendrautumėte su kitu žaidėju dėl savo tikslinės būsenos! Vienintelis būdas, kuriuo galite perduoti savo strategiją kitam žaidėjui, yra komanda `SAKYTI: <ŽINUTĖ>`!} \\
\\ 
\\ 
\texttt{**Objektų perkėlimas**} \\
\\ 
\texttt{* Objektus galima perkelti tik į langelius, esančius tinklelio ribose. Tikslinis langelis turi būti tuščias, t. y. jame turi būti tik simbolis '◌'.} \\
\texttt{* Jei bandysite perkelti objektą į vietą, kuri nėra tuščia, arba bandysite jį perkelti už tinklelio ribų, turėsiu jus nubausti. Gausite dar vieną bandymą.} \\
\texttt{* Prieš atlikdami perkėlimą, dukart patikrinkite, ar tikslinė vieta yra tuščia ir ar joje nėra jokios raidės, rėmelio ar linijos!} \\
\\ 
\\ 
\texttt{**Žaidimo pabaiga**} \\
\\ 
\texttt{Jei manote, kad pasiekėte tikslą {-} sulygiavote visus objektus, galite paprašyti kito žaidėjo baigti žaidimą siųsdami `SAKYTI: baigta?`. Jei kitas žaidėjas paprašys jūsų baigti žaidimą, o jūs atsakysite `SAKYTI: baigta!`, žaidimas baigsis.} \\
\\ 
\texttt{Abu žaidėjai laimi, jei žaidimas baigiasi per 20 raundus, kur vienas raundas apibrėžiamas kaip du žaidėjai, kiekvienas išsiuntęs po vieną galiojančią komandą.} \\
\\ 
\texttt{**Taškų skaičiavimas:**} \\
\\ 
\texttt{Kuo arčiau vienodų objektų abiejose žaidimo lentose, tuo daugiau taškų gausite. Baudos mažina jūsų taškus. Ar pavyks pagerinti rekordą?} \\
\\ 
\texttt{Kitas žaidėjas pradėjo žaidimą išsiųsdamas šią žinutę:} \\
\\ 
\texttt{"Sveiki! Kadangi mūsų lentos skiriasi, mums reikia susitarti dėl bendrų tikslinių pozicijų. Prašau, pasiūlykite 5 konkrečias (X, Y) koordinates, kuriose jūsų lentos yra laisvi langeliai ('◌'). Aš tuos pačius koordinates naudosiu savo lentos laisvoms vietoms perkelti savo objektus. Tada mūsų objektai sutaps. Laukiu jūsų koordinatžių!"} \\
\\ 
\texttt{Kokia yra jūsų pirmoji komanda?} \\
            }
        }
    }
    & \\ \\

    \theutterance \stepcounter{utterance}  
    & & & \multicolumn{4}{p{0.6\linewidth}}{
        \cellcolor[rgb]{0.9,0.9,0.9}{
            \makecell[{{p{\linewidth}}}]{
                \texttt{\tiny{[GM$\langle$P2]}}
                \texttt{SAKYTI: Sveiki! Sutinku. Štai 5 koordinatės, kurias siūlau naudoti kaip bendras tikslines pozicijas mūsų lentose: (2, 1), (3, 2), (4, 3), (5, 4), (6, 5). Prašau patvirtinti, ar šios koordinatės tinka ir ar jūsų lentose šiose vietose yra laisvų langelių ('◌').} \\
            }
        }
    }
     \\ \\

    \theutterance \stepcounter{utterance}  
    & & & \multicolumn{2}{p{0.3\linewidth}}{
        \cellcolor[rgb]{0.9,0.9,0.9}{
            \makecell[{{p{\linewidth}}}]{
                \texttt{\tiny{[GM$|$GM]}}
                \texttt{SAKYTI: Sveiki! Sutinku. Štai 5 koordinatės, kurias siūlau naudoti kaip bendras tikslines pozicijas mūsų lentose: (2, 1), (3, 2), (4, 3), (5, 4), (6, 5). Prašau patvirtinti, ar šios koordinatės tinka ir ar jūsų lentose šiose vietose yra laisvų langelių ('◌').} \\
            }
        }
    }
    & & \\ \\

    \theutterance \stepcounter{utterance}  
    & & \multicolumn{4}{p{0.6\linewidth}}{
        \cellcolor[rgb]{0.9,0.9,0.9}{
            \makecell[{{p{\linewidth}}}]{
                \texttt{\tiny{[P1$\langle$GM]}}
                \texttt{Jūsų žinutė buvo perduota kitam žaidėjui.} \\
\\ 
\\ 
\texttt{Šiuo metu žaidžiate 1 raundą iš maksimaliai 20.} \\
\texttt{Jūs bendrai sukaupėte 1 iš 12 baudų.} \\
\texttt{Kitas žaidėjas atsiuntė šią žinutę:} \\
\\ 
\\ 
\texttt{"Sveiki! Sutinku. Štai 5 koordinatės, kurias siūlau naudoti kaip bendras tikslines pozicijas mūsų lentose: (2, 1), (3, 2), (4, 3), (5, 4), (6, 5). Prašau patvirtinti, ar šios koordinatės tinka ir ar jūsų lentose šiose vietose yra laisvų langelių ('◌')."} \\
\\ 
\\ 
\texttt{Kokia jūsų kita komanda?} \\
            }
        }
    }
    & \\ \\

    \theutterance \stepcounter{utterance}  
    & \multicolumn{4}{p{0.6\linewidth}}{
        \cellcolor[rgb]{0.9,0.9,0.9}{
            \makecell[{{p{\linewidth}}}]{
                \texttt{\tiny{[P1$\rangle$GM]}}
                \texttt{SAKYTI: Patvirtinu! Šios koordinatės tinka. Patikrinau savo lentą:} \\
\texttt{(2, 1) yra laisva.} \\
\texttt{(3, 2) yra laisva.} \\
\texttt{(4, 3) yra laisva.} \\
\texttt{(5, 4) yra laisva.} \\
\texttt{(6, 5) yra laisva.} \\
\\ 
\texttt{Pradedame perkelti objektus į šias pozicijas. Mano planas:} \\
\texttt{1. Perkelti 'C' į (2, 1)} \\
\texttt{2. Perkelti 'I' į (3, 2)} \\
\texttt{3. Perkelti 'T' į (4, 3)} \\
\texttt{4. Perkelti 'W' į (5, 4)} \\
\texttt{5. Perkelti 'H' į (6, 5)} \\
\\ 
\texttt{Pirmiausia perkelsiu 'C' į (2, 1).} \\
            }
        }
    }
    & & \\ \\

    \theutterance \stepcounter{utterance}  
    & & & \multicolumn{2}{p{0.3\linewidth}}{
        \cellcolor[rgb]{0.9,0.9,0.9}{
            \makecell[{{p{\linewidth}}}]{
                \texttt{\tiny{[GM$|$GM]}}
                \texttt{SAKYTI: Patvirtinu! Šios koordinatės tinka. Patikrinau savo lentą:} \\
\texttt{(2, 1) yra laisva.} \\
\texttt{(3, 2) yra laisva.} \\
\texttt{(4, 3) yra laisva.} \\
\texttt{(5, 4) yra laisva.} \\
\texttt{(6, 5) yra laisva.} \\
\\ 
\texttt{Pradedame perkelti objektus į šias pozicijas. Mano planas:} \\
\texttt{1. Perkelti 'C' į (2, 1)} \\
\texttt{2. Perkelti 'I' į (3, 2)} \\
\texttt{3. Perkelti 'T' į (4, 3)} \\
\texttt{4. Perkelti 'W' į (5, 4)} \\
\texttt{5. Perkelti 'H' į (6, 5)} \\
\\ 
\texttt{Pirmiausia perkelsiu 'C' į (2, 1).} \\
            }
        }
    }
    & & \\ \\

    \theutterance \stepcounter{utterance}  
    & & & \multicolumn{2}{p{0.3\linewidth}}{
        \cellcolor[rgb]{0.9,0.9,0.9}{
            \makecell[{{p{\linewidth}}}]{
                \texttt{\tiny{[GM$|$GM]}}
                \texttt{Invalid format: tail is not empty: '} \\
\texttt{(2, 1) yra laisva.} \\
\texttt{(3, 2) yra laisva.} \\
\texttt{(4, 3) yra laisva.} \\
\texttt{(5, 4) yra laisva.} \\
\texttt{(6, 5) yra laisva.} \\
\\ 
\texttt{Pradedame perkelti objektus į šias pozicijas. Mano planas:} \\
\texttt{1. Perkelti 'C' į (2, 1)} \\
\texttt{2. Perkelti 'I' į (3, 2)} \\
\texttt{3. Perkelti 'T' į (4, 3)} \\
\texttt{4. Perkelti 'W' į (5, 4)} \\
\texttt{5. Perkelti 'H' į (6, 5)} \\
\\ 
\texttt{Pirmiausia perkelsiu 'C' į (2, 1).} \\
            }
        }
    }
    & & \\ \\

    \theutterance \stepcounter{utterance}  
    & & \multicolumn{4}{p{0.6\linewidth}}{
        \cellcolor[rgb]{0.9,0.9,0.9}{
            \makecell[{{p{\linewidth}}}]{
                \texttt{\tiny{[P1$\langle$GM]}}
                \texttt{Nuobauda: Jūsų žinutėje neturi būti nieko po komandos!} \\
\\ 
\texttt{Įsitikinkite, kad jūsų atsakyme yra tik `SAKYTI: <ŽINUTĖ>` arba `JUDĖTI: <OBJEKTAS>, (<X>, <Y>)`, ir nieko daugiau!} \\
\\ 
\texttt{Jūs bendrai sukaupėte 2 iš 12 baudų.} \\
\texttt{Pabandykite dar kartą!} \\
            }
        }
    }
    & \\ \\

    \theutterance \stepcounter{utterance}  
    & \multicolumn{4}{p{0.6\linewidth}}{
        \cellcolor[rgb]{0.9,0.9,0.9}{
            \makecell[{{p{\linewidth}}}]{
                \texttt{\tiny{[P1$\rangle$GM]}}
                \texttt{JUDĖTI: C, (2, 1)} \\
            }
        }
    }
    & & \\ \\

    \theutterance \stepcounter{utterance}  
    & & & \multicolumn{2}{p{0.3\linewidth}}{
        \cellcolor[rgb]{0.9,0.9,0.9}{
            \makecell[{{p{\linewidth}}}]{
                \texttt{\tiny{[GM$|$GM]}}
                \texttt{JUDĖTI: C, (2, 1)} \\
            }
        }
    }
    & & \\ \\

    \theutterance \stepcounter{utterance}  
    & & & \multicolumn{2}{p{0.3\linewidth}}{
        \cellcolor[rgb]{0.9,0.9,0.9}{
            \makecell[{{p{\linewidth}}}]{
                \texttt{\tiny{[GM$|$GM]}}
                \texttt{C moved to (2, 1)} \\
            }
        }
    }
    & & \\ \\

    \theutterance \stepcounter{utterance}  
    & & \multicolumn{4}{p{0.6\linewidth}}{
        \cellcolor[rgb]{0.9,0.9,0.9}{
            \makecell[{{p{\linewidth}}}]{
                \texttt{\tiny{[GM$\rangle$P2]}}
                \texttt{Jūsų žinutė buvo perduota kitam žaidėjui.} \\
\\ 
\\ 
\texttt{Šiuo metu žaidžiate 2 raundą iš maksimaliai 20.} \\
\texttt{Jūs bendrai sukaupėte 2 iš 12 baudų.} \\
\texttt{Kitas žaidėjas perkėlė objektą savo tinklelyje.} \\
\\ 
\\ 
\texttt{Kokia jūsų kita komanda?} \\
            }
        }
    }
    & \\ \\

    \theutterance \stepcounter{utterance}  
    & & & \multicolumn{4}{p{0.6\linewidth}}{
        \cellcolor[rgb]{0.9,0.9,0.9}{
            \makecell[{{p{\linewidth}}}]{
                \texttt{\tiny{[GM$\langle$P2]}}
                \texttt{SAKYTI: Koordinatės (2, 1), (3, 2), (4, 3), (5, 4), (6, 5) man tinka. Prašau patvirtinti, kad ir jūsų lentose šiose vietose yra laisvų langelių. Jei taip, galime pradėti perkelti objektus į šias pozicijas. Pirmiausia, kurį objektą norėtumėte perkelti pirmąjį?} \\
            }
        }
    }
     \\ \\

    \theutterance \stepcounter{utterance}  
    & & & \multicolumn{2}{p{0.3\linewidth}}{
        \cellcolor[rgb]{0.9,0.9,0.9}{
            \makecell[{{p{\linewidth}}}]{
                \texttt{\tiny{[GM$|$GM]}}
                \texttt{SAKYTI: Koordinatės (2, 1), (3, 2), (4, 3), (5, 4), (6, 5) man tinka. Prašau patvirtinti, kad ir jūsų lentose šiose vietose yra laisvų langelių. Jei taip, galime pradėti perkelti objektus į šias pozicijas. Pirmiausia, kurį objektą norėtumėte perkelti pirmąjį?} \\
            }
        }
    }
    & & \\ \\

    \theutterance \stepcounter{utterance}  
    & & \multicolumn{4}{p{0.6\linewidth}}{
        \cellcolor[rgb]{0.9,0.9,0.9}{
            \makecell[{{p{\linewidth}}}]{
                \texttt{\tiny{[P1$\langle$GM]}}
                \texttt{Sėkmingai perkelta 'C' į (2,1). Jūsų atnaujintas tinklelis atrodo taip:} \\
\texttt{```} \\
\texttt{ 1234567} \\
\texttt{╔══╤════╗} \\
\texttt{║◌C│◌◌◌H║ 1} \\
\texttt{╟──┤◌◌┌─╢ 2} \\
\texttt{║◌I├──┤◌║ 3} \\
\texttt{╟──┘◌◌├─╢ 4} \\
\texttt{║◌◌◌◌◌│◌║ 5} \\
\texttt{╟─┐W┌─┼─╢ 6} \\
\texttt{║◌│T│◌│◌║ 7} \\
\texttt{╚═╧═╧═╧═╝} \\
\texttt{```} \\
\texttt{Šiuo metu žaidžiate 2 raundą iš maksimaliai 20.} \\
\texttt{Jūs bendrai sukaupėte 2 iš 12 baudų.} \\
\texttt{Kitas žaidėjas atsiuntė šią žinutę:} \\
\\ 
\\ 
\texttt{"Koordinatės (2, 1), (3, 2), (4, 3), (5, 4), (6, 5) man tinka. Prašau patvirtinti, kad ir jūsų lentose šiose vietose yra laisvų langelių. Jei taip, galime pradėti perkelti objektus į šias pozicijas. Pirmiausia, kurį objektą norėtumėte perkelti pirmąjį?"} \\
\\ 
\\ 
\texttt{Kokia jūsų kita komanda?} \\
            }
        }
    }
    & \\ \\

    \theutterance \stepcounter{utterance}  
    & \multicolumn{4}{p{0.6\linewidth}}{
        \cellcolor[rgb]{0.9,0.9,0.9}{
            \makecell[{{p{\linewidth}}}]{
                \texttt{\tiny{[P1$\rangle$GM]}}
                \texttt{SAKYTI: Patvirtinu, kad koordinatės tinka. Aš jau perkėliau 'C' į (2, 1). Dabar perkelsiu 'I' į (3, 2).} \\
            }
        }
    }
    & & \\ \\

    \theutterance \stepcounter{utterance}  
    & & & \multicolumn{2}{p{0.3\linewidth}}{
        \cellcolor[rgb]{0.9,0.9,0.9}{
            \makecell[{{p{\linewidth}}}]{
                \texttt{\tiny{[GM$|$GM]}}
                \texttt{SAKYTI: Patvirtinu, kad koordinatės tinka. Aš jau perkėliau 'C' į (2, 1). Dabar perkelsiu 'I' į (3, 2).} \\
            }
        }
    }
    & & \\ \\

    \theutterance \stepcounter{utterance}  
    & & \multicolumn{4}{p{0.6\linewidth}}{
        \cellcolor[rgb]{0.9,0.9,0.9}{
            \makecell[{{p{\linewidth}}}]{
                \texttt{\tiny{[GM$\rangle$P2]}}
                \texttt{Jūsų žinutė buvo perduota kitam žaidėjui.} \\
\\ 
\\ 
\texttt{Šiuo metu žaidžiate 3 raundą iš maksimaliai 20.} \\
\texttt{Jūs bendrai sukaupėte 2 iš 12 baudų.} \\
\texttt{Kitas žaidėjas atsiuntė šią žinutę:} \\
\\ 
\\ 
\texttt{"Patvirtinu, kad koordinatės tinka. Aš jau perkėliau 'C' į (2, 1). Dabar perkelsiu 'I' į (3, 2)."} \\
\\ 
\\ 
\texttt{Kokia jūsų kita komanda?} \\
            }
        }
    }
    & \\ \\

    \theutterance \stepcounter{utterance}  
    & & & \multicolumn{4}{p{0.6\linewidth}}{
        \cellcolor[rgb]{0.9,0.9,0.9}{
            \makecell[{{p{\linewidth}}}]{
                \texttt{\tiny{[GM$\langle$P2]}}
                \texttt{SAKYTI: Puiku! Kadangi jūs jau perkėlėte 'C' į (2, 1) ir planuojate perkelti 'I' į (3, 2), aš taip pat atliksiu tuos pačius veiksmus savo lentelėje. Pirmiausia perkelsiu objektą 'C' į poziciją (2, 1).} \\
            }
        }
    }
     \\ \\

    \theutterance \stepcounter{utterance}  
    & & & \multicolumn{2}{p{0.3\linewidth}}{
        \cellcolor[rgb]{0.9,0.9,0.9}{
            \makecell[{{p{\linewidth}}}]{
                \texttt{\tiny{[GM$|$GM]}}
                \texttt{SAKYTI: Puiku! Kadangi jūs jau perkėlėte 'C' į (2, 1) ir planuojate perkelti 'I' į (3, 2), aš taip pat atliksiu tuos pačius veiksmus savo lentelėje. Pirmiausia perkelsiu objektą 'C' į poziciją (2, 1).} \\
            }
        }
    }
    & & \\ \\

    \theutterance \stepcounter{utterance}  
    & & \multicolumn{4}{p{0.6\linewidth}}{
        \cellcolor[rgb]{0.9,0.9,0.9}{
            \makecell[{{p{\linewidth}}}]{
                \texttt{\tiny{[P1$\langle$GM]}}
                \texttt{Jūsų žinutė buvo perduota kitam žaidėjui.} \\
\\ 
\\ 
\texttt{Šiuo metu žaidžiate 3 raundą iš maksimaliai 20.} \\
\texttt{Jūs bendrai sukaupėte 2 iš 12 baudų.} \\
\texttt{Kitas žaidėjas atsiuntė šią žinutę:} \\
\\ 
\\ 
\texttt{"Puiku! Kadangi jūs jau perkėlėte 'C' į (2, 1) ir planuojate perkelti 'I' į (3, 2), aš taip pat atliksiu tuos pačius veiksmus savo lentelėje. Pirmiausia perkelsiu objektą 'C' į poziciją (2, 1)."} \\
\\ 
\\ 
\texttt{Kokia jūsų kita komanda?} \\
            }
        }
    }
    & \\ \\

    \theutterance \stepcounter{utterance}  
    & \multicolumn{4}{p{0.6\linewidth}}{
        \cellcolor[rgb]{0.9,0.9,0.9}{
            \makecell[{{p{\linewidth}}}]{
                \texttt{\tiny{[P1$\rangle$GM]}}
                \texttt{JUDĖTI: I, (3, 2)} \\
            }
        }
    }
    & & \\ \\

    \theutterance \stepcounter{utterance}  
    & & & \multicolumn{2}{p{0.3\linewidth}}{
        \cellcolor[rgb]{0.9,0.9,0.9}{
            \makecell[{{p{\linewidth}}}]{
                \texttt{\tiny{[GM$|$GM]}}
                \texttt{JUDĖTI: I, (3, 2)} \\
            }
        }
    }
    & & \\ \\

    \theutterance \stepcounter{utterance}  
    & & & \multicolumn{2}{p{0.3\linewidth}}{
        \cellcolor[rgb]{0.9,0.9,0.9}{
            \makecell[{{p{\linewidth}}}]{
                \texttt{\tiny{[GM$|$GM]}}
                \texttt{Nuobauda: (3,2) nėra tuščias, bet jame yra '┤'.} \\
\\ 
\texttt{Jūs bendrai sukaupėte 3 iš 12 baudų.} \\
\texttt{Pabandykite dar kartą!} \\
            }
        }
    }
    & & \\ \\

    \theutterance \stepcounter{utterance}  
    & & \multicolumn{4}{p{0.6\linewidth}}{
        \cellcolor[rgb]{0.9,0.9,0.9}{
            \makecell[{{p{\linewidth}}}]{
                \texttt{\tiny{[P1$\langle$GM]}}
                \texttt{Nuobauda: (3,2) nėra tuščias, bet jame yra '┤'.} \\
\\ 
\texttt{Jūs bendrai sukaupėte 3 iš 12 baudų.} \\
\texttt{Pabandykite dar kartą!} \\
            }
        }
    }
    & \\ \\

    \theutterance \stepcounter{utterance}  
    & \multicolumn{4}{p{0.6\linewidth}}{
        \cellcolor[rgb]{0.9,0.9,0.9}{
            \makecell[{{p{\linewidth}}}]{
                \texttt{\tiny{[P1$\rangle$GM]}}
                \texttt{JUDĖTI: I, (3, 1)} \\
            }
        }
    }
    & & \\ \\

    \theutterance \stepcounter{utterance}  
    & & & \multicolumn{2}{p{0.3\linewidth}}{
        \cellcolor[rgb]{0.9,0.9,0.9}{
            \makecell[{{p{\linewidth}}}]{
                \texttt{\tiny{[GM$|$GM]}}
                \texttt{JUDĖTI: I, (3, 1)} \\
            }
        }
    }
    & & \\ \\

    \theutterance \stepcounter{utterance}  
    & & & \multicolumn{2}{p{0.3\linewidth}}{
        \cellcolor[rgb]{0.9,0.9,0.9}{
            \makecell[{{p{\linewidth}}}]{
                \texttt{\tiny{[GM$|$GM]}}
                \texttt{Nuobauda: (3,1) nėra tuščias, bet jame yra '│'.} \\
\\ 
\texttt{Jūs bendrai sukaupėte 4 iš 12 baudų.} \\
\texttt{Pabandykite dar kartą!} \\
            }
        }
    }
    & & \\ \\

    \theutterance \stepcounter{utterance}  
    & & \multicolumn{4}{p{0.6\linewidth}}{
        \cellcolor[rgb]{0.9,0.9,0.9}{
            \makecell[{{p{\linewidth}}}]{
                \texttt{\tiny{[P1$\langle$GM]}}
                \texttt{Nuobauda: (3,1) nėra tuščias, bet jame yra '│'.} \\
\\ 
\texttt{Jūs bendrai sukaupėte 4 iš 12 baudų.} \\
\texttt{Pabandykite dar kartą!} \\
            }
        }
    }
    & \\ \\

    \theutterance \stepcounter{utterance}  
    & \multicolumn{4}{p{0.6\linewidth}}{
        \cellcolor[rgb]{0.9,0.9,0.9}{
            \makecell[{{p{\linewidth}}}]{
                \texttt{\tiny{[P1$\rangle$GM]}}
                \texttt{JUDĖTI: I, (2, 2)} \\
            }
        }
    }
    & & \\ \\

    \theutterance \stepcounter{utterance}  
    & & & \multicolumn{2}{p{0.3\linewidth}}{
        \cellcolor[rgb]{0.9,0.9,0.9}{
            \makecell[{{p{\linewidth}}}]{
                \texttt{\tiny{[GM$|$GM]}}
                \texttt{JUDĖTI: I, (2, 2)} \\
            }
        }
    }
    & & \\ \\

    \theutterance \stepcounter{utterance}  
    & & & \multicolumn{2}{p{0.3\linewidth}}{
        \cellcolor[rgb]{0.9,0.9,0.9}{
            \makecell[{{p{\linewidth}}}]{
                \texttt{\tiny{[GM$|$GM]}}
                \texttt{Nuobauda: (2,2) nėra tuščias, bet jame yra '─'.} \\
\\ 
\texttt{Jūs bendrai sukaupėte 5 iš 12 baudų.} \\
\texttt{Pabandykite dar kartą!} \\
            }
        }
    }
    & & \\ \\

    \theutterance \stepcounter{utterance}  
    & & \multicolumn{4}{p{0.6\linewidth}}{
        \cellcolor[rgb]{0.9,0.9,0.9}{
            \makecell[{{p{\linewidth}}}]{
                \texttt{\tiny{[P1$\langle$GM]}}
                \texttt{Nuobauda: (2,2) nėra tuščias, bet jame yra '─'.} \\
\\ 
\texttt{Jūs bendrai sukaupėte 5 iš 12 baudų.} \\
\texttt{Pabandykite dar kartą!} \\
            }
        }
    }
    & \\ \\

    \theutterance \stepcounter{utterance}  
    & \multicolumn{4}{p{0.6\linewidth}}{
        \cellcolor[rgb]{0.9,0.9,0.9}{
            \makecell[{{p{\linewidth}}}]{
                \texttt{\tiny{[P1$\rangle$GM]}}
                \texttt{JUDĖTI: I, (3, 3)} \\
            }
        }
    }
    & & \\ \\

    \theutterance \stepcounter{utterance}  
    & & & \multicolumn{2}{p{0.3\linewidth}}{
        \cellcolor[rgb]{0.9,0.9,0.9}{
            \makecell[{{p{\linewidth}}}]{
                \texttt{\tiny{[GM$|$GM]}}
                \texttt{JUDĖTI: I, (3, 3)} \\
            }
        }
    }
    & & \\ \\

    \theutterance \stepcounter{utterance}  
    & & & \multicolumn{2}{p{0.3\linewidth}}{
        \cellcolor[rgb]{0.9,0.9,0.9}{
            \makecell[{{p{\linewidth}}}]{
                \texttt{\tiny{[GM$|$GM]}}
                \texttt{Nuobauda: (3,3) nėra tuščias, bet jame yra '├'.} \\
\\ 
\texttt{Jūs bendrai sukaupėte 6 iš 12 baudų.} \\
\texttt{Pabandykite dar kartą!} \\
            }
        }
    }
    & & \\ \\

    \theutterance \stepcounter{utterance}  
    & & \multicolumn{4}{p{0.6\linewidth}}{
        \cellcolor[rgb]{0.9,0.9,0.9}{
            \makecell[{{p{\linewidth}}}]{
                \texttt{\tiny{[P1$\langle$GM]}}
                \texttt{Nuobauda: (3,3) nėra tuščias, bet jame yra '├'.} \\
\\ 
\texttt{Jūs bendrai sukaupėte 6 iš 12 baudų.} \\
\texttt{Pabandykite dar kartą!} \\
            }
        }
    }
    & \\ \\

    \theutterance \stepcounter{utterance}  
    & \multicolumn{4}{p{0.6\linewidth}}{
        \cellcolor[rgb]{0.9,0.9,0.9}{
            \makecell[{{p{\linewidth}}}]{
                \texttt{\tiny{[P1$\rangle$GM]}}
                \texttt{JUDĖTI: I, (4, 2)} \\
            }
        }
    }
    & & \\ \\

    \theutterance \stepcounter{utterance}  
    & & & \multicolumn{2}{p{0.3\linewidth}}{
        \cellcolor[rgb]{0.9,0.9,0.9}{
            \makecell[{{p{\linewidth}}}]{
                \texttt{\tiny{[GM$|$GM]}}
                \texttt{JUDĖTI: I, (4, 2)} \\
            }
        }
    }
    & & \\ \\

    \theutterance \stepcounter{utterance}  
    & & & \multicolumn{2}{p{0.3\linewidth}}{
        \cellcolor[rgb]{0.9,0.9,0.9}{
            \makecell[{{p{\linewidth}}}]{
                \texttt{\tiny{[GM$|$GM]}}
                \texttt{I moved to (4, 2)} \\
            }
        }
    }
    & & \\ \\

    \theutterance \stepcounter{utterance}  
    & & \multicolumn{4}{p{0.6\linewidth}}{
        \cellcolor[rgb]{0.9,0.9,0.9}{
            \makecell[{{p{\linewidth}}}]{
                \texttt{\tiny{[GM$\rangle$P2]}}
                \texttt{Jūsų žinutė buvo perduota kitam žaidėjui.} \\
\\ 
\\ 
\texttt{Šiuo metu žaidžiate 4 raundą iš maksimaliai 20.} \\
\texttt{Jūs bendrai sukaupėte 6 iš 12 baudų.} \\
\texttt{Kitas žaidėjas perkėlė objektą savo tinklelyje.} \\
\\ 
\\ 
\texttt{Kokia jūsų kita komanda?} \\
            }
        }
    }
    & \\ \\

    \theutterance \stepcounter{utterance}  
    & & & \multicolumn{4}{p{0.6\linewidth}}{
        \cellcolor[rgb]{0.9,0.9,0.9}{
            \makecell[{{p{\linewidth}}}]{
                \texttt{\tiny{[GM$\langle$P2]}}
                \texttt{JUDĖTI: C, (2, 1)} \\
            }
        }
    }
     \\ \\

    \theutterance \stepcounter{utterance}  
    & & & \multicolumn{2}{p{0.3\linewidth}}{
        \cellcolor[rgb]{0.9,0.9,0.9}{
            \makecell[{{p{\linewidth}}}]{
                \texttt{\tiny{[GM$|$GM]}}
                \texttt{JUDĖTI: C, (2, 1)} \\
            }
        }
    }
    & & \\ \\

    \theutterance \stepcounter{utterance}  
    & & & \multicolumn{2}{p{0.3\linewidth}}{
        \cellcolor[rgb]{0.9,0.9,0.9}{
            \makecell[{{p{\linewidth}}}]{
                \texttt{\tiny{[GM$|$GM]}}
                \texttt{Nuobauda: (2,1) nėra tuščias, bet jame yra 'W'.} \\
\\ 
\texttt{Jūs bendrai sukaupėte 7 iš 12 baudų.} \\
\texttt{Pabandykite dar kartą!} \\
            }
        }
    }
    & & \\ \\

    \theutterance \stepcounter{utterance}  
    & & \multicolumn{4}{p{0.6\linewidth}}{
        \cellcolor[rgb]{0.9,0.9,0.9}{
            \makecell[{{p{\linewidth}}}]{
                \texttt{\tiny{[GM$\rangle$P2]}}
                \texttt{Nuobauda: (2,1) nėra tuščias, bet jame yra 'W'.} \\
\\ 
\texttt{Jūs bendrai sukaupėte 7 iš 12 baudų.} \\
\texttt{Pabandykite dar kartą!} \\
            }
        }
    }
    & \\ \\

    \theutterance \stepcounter{utterance}  
    & & & \multicolumn{4}{p{0.6\linewidth}}{
        \cellcolor[rgb]{0.9,0.9,0.9}{
            \makecell[{{p{\linewidth}}}]{
                \texttt{\tiny{[GM$\langle$P2]}}
                \texttt{JUDĖTI: C, (3, 2)} \\
            }
        }
    }
     \\ \\

    \theutterance \stepcounter{utterance}  
    & & & \multicolumn{2}{p{0.3\linewidth}}{
        \cellcolor[rgb]{0.9,0.9,0.9}{
            \makecell[{{p{\linewidth}}}]{
                \texttt{\tiny{[GM$|$GM]}}
                \texttt{JUDĖTI: C, (3, 2)} \\
            }
        }
    }
    & & \\ \\

    \theutterance \stepcounter{utterance}  
    & & & \multicolumn{2}{p{0.3\linewidth}}{
        \cellcolor[rgb]{0.9,0.9,0.9}{
            \makecell[{{p{\linewidth}}}]{
                \texttt{\tiny{[GM$|$GM]}}
                \texttt{Nuobauda: (3,2) nėra tuščias, bet jame yra '┤'.} \\
\\ 
\texttt{Jūs bendrai sukaupėte 8 iš 12 baudų.} \\
\texttt{Pabandykite dar kartą!} \\
            }
        }
    }
    & & \\ \\

    \theutterance \stepcounter{utterance}  
    & & \multicolumn{4}{p{0.6\linewidth}}{
        \cellcolor[rgb]{0.9,0.9,0.9}{
            \makecell[{{p{\linewidth}}}]{
                \texttt{\tiny{[GM$\rangle$P2]}}
                \texttt{Nuobauda: (3,2) nėra tuščias, bet jame yra '┤'.} \\
\\ 
\texttt{Jūs bendrai sukaupėte 8 iš 12 baudų.} \\
\texttt{Pabandykite dar kartą!} \\
            }
        }
    }
    & \\ \\

    \theutterance \stepcounter{utterance}  
    & & & \multicolumn{4}{p{0.6\linewidth}}{
        \cellcolor[rgb]{0.9,0.9,0.9}{
            \makecell[{{p{\linewidth}}}]{
                \texttt{\tiny{[GM$\langle$P2]}}
                \texttt{JUDĖTI: C, (4, 3)} \\
            }
        }
    }
     \\ \\

    \theutterance \stepcounter{utterance}  
    & & & \multicolumn{2}{p{0.3\linewidth}}{
        \cellcolor[rgb]{0.9,0.9,0.9}{
            \makecell[{{p{\linewidth}}}]{
                \texttt{\tiny{[GM$|$GM]}}
                \texttt{JUDĖTI: C, (4, 3)} \\
            }
        }
    }
    & & \\ \\

    \theutterance \stepcounter{utterance}  
    & & & \multicolumn{2}{p{0.3\linewidth}}{
        \cellcolor[rgb]{0.9,0.9,0.9}{
            \makecell[{{p{\linewidth}}}]{
                \texttt{\tiny{[GM$|$GM]}}
                \texttt{Nuobauda: (4,3) nėra tuščias, bet jame yra '─'.} \\
\\ 
\texttt{Jūs bendrai sukaupėte 9 iš 12 baudų.} \\
\texttt{Pabandykite dar kartą!} \\
            }
        }
    }
    & & \\ \\

    \theutterance \stepcounter{utterance}  
    & & \multicolumn{4}{p{0.6\linewidth}}{
        \cellcolor[rgb]{0.9,0.9,0.9}{
            \makecell[{{p{\linewidth}}}]{
                \texttt{\tiny{[GM$\rangle$P2]}}
                \texttt{Nuobauda: (4,3) nėra tuščias, bet jame yra '─'.} \\
\\ 
\texttt{Jūs bendrai sukaupėte 9 iš 12 baudų.} \\
\texttt{Pabandykite dar kartą!} \\
            }
        }
    }
    & \\ \\

    \theutterance \stepcounter{utterance}  
    & & & \multicolumn{4}{p{0.6\linewidth}}{
        \cellcolor[rgb]{0.9,0.9,0.9}{
            \makecell[{{p{\linewidth}}}]{
                \texttt{\tiny{[GM$\langle$P2]}}
                \texttt{JUDĖTI: C, (5, 4)} \\
            }
        }
    }
     \\ \\

    \theutterance \stepcounter{utterance}  
    & & & \multicolumn{2}{p{0.3\linewidth}}{
        \cellcolor[rgb]{0.9,0.9,0.9}{
            \makecell[{{p{\linewidth}}}]{
                \texttt{\tiny{[GM$|$GM]}}
                \texttt{JUDĖTI: C, (5, 4)} \\
            }
        }
    }
    & & \\ \\

    \theutterance \stepcounter{utterance}  
    & & & \multicolumn{2}{p{0.3\linewidth}}{
        \cellcolor[rgb]{0.9,0.9,0.9}{
            \makecell[{{p{\linewidth}}}]{
                \texttt{\tiny{[GM$|$GM]}}
                \texttt{C moved to (5, 4)} \\
            }
        }
    }
    & & \\ \\

    \theutterance \stepcounter{utterance}  
    & & \multicolumn{4}{p{0.6\linewidth}}{
        \cellcolor[rgb]{0.9,0.9,0.9}{
            \makecell[{{p{\linewidth}}}]{
                \texttt{\tiny{[P1$\langle$GM]}}
                \texttt{Sėkmingai perkelta 'I' į (4,2). Jūsų atnaujintas tinklelis atrodo taip:} \\
\texttt{```} \\
\texttt{ 1234567} \\
\texttt{╔══╤════╗} \\
\texttt{║◌C│◌◌◌H║ 1} \\
\texttt{╟──┤I◌┌─╢ 2} \\
\texttt{║◌◌├──┤◌║ 3} \\
\texttt{╟──┘◌◌├─╢ 4} \\
\texttt{║◌◌◌◌◌│◌║ 5} \\
\texttt{╟─┐W┌─┼─╢ 6} \\
\texttt{║◌│T│◌│◌║ 7} \\
\texttt{╚═╧═╧═╧═╝} \\
\texttt{```} \\
\texttt{Šiuo metu žaidžiate 4 raundą iš maksimaliai 20.} \\
\texttt{Jūs bendrai sukaupėte 9 iš 12 baudų.} \\
\texttt{Kitas žaidėjas perkėlė objektą savo tinklelyje.} \\
\\ 
\\ 
\texttt{Kokia jūsų kita komanda?} \\
            }
        }
    }
    & \\ \\

    \theutterance \stepcounter{utterance}  
    & \multicolumn{4}{p{0.6\linewidth}}{
        \cellcolor[rgb]{0.9,0.9,0.9}{
            \makecell[{{p{\linewidth}}}]{
                \texttt{\tiny{[P1$\rangle$GM]}}
                \texttt{SAKYTI: Dabar perkelsiu 'T' į (4, 3). Patikrinkite, ar (4,3) yra laisva vieta jūsų lentelėje.} \\
            }
        }
    }
    & & \\ \\

    \theutterance \stepcounter{utterance}  
    & & & \multicolumn{2}{p{0.3\linewidth}}{
        \cellcolor[rgb]{0.9,0.9,0.9}{
            \makecell[{{p{\linewidth}}}]{
                \texttt{\tiny{[GM$|$GM]}}
                \texttt{SAKYTI: Dabar perkelsiu 'T' į (4, 3). Patikrinkite, ar (4,3) yra laisva vieta jūsų lentelėje.} \\
            }
        }
    }
    & & \\ \\

    \theutterance \stepcounter{utterance}  
    & & \multicolumn{4}{p{0.6\linewidth}}{
        \cellcolor[rgb]{0.9,0.9,0.9}{
            \makecell[{{p{\linewidth}}}]{
                \texttt{\tiny{[GM$\rangle$P2]}}
                \texttt{Sėkmingai perkelta 'C' į (5,4). Jūsų atnaujintas tinklelis atrodo taip:} \\
\texttt{```} \\
\texttt{ 1234567} \\
\texttt{╔══╤════╗} \\
\texttt{║◌W│◌◌◌◌║ 1} \\
\texttt{╟──┤◌H┌─╢ 2} \\
\texttt{║◌◌├──┤◌║ 3} \\
\texttt{╟──┘◌C├─╢ 4} \\
\texttt{║◌◌◌◌◌│◌║ 5} \\
\texttt{╟─┐◌┌─┼─╢ 6} \\
\texttt{║I│◌│T│◌║ 7} \\
\texttt{╚═╧═╧═╧═╝} \\
\texttt{```} \\
\texttt{Šiuo metu žaidžiate 5 raundą iš maksimaliai 20.} \\
\texttt{Jūs bendrai sukaupėte 9 iš 12 baudų.} \\
\texttt{Kitas žaidėjas atsiuntė šią žinutę:} \\
\\ 
\\ 
\texttt{"Dabar perkelsiu 'T' į (4, 3). Patikrinkite, ar (4,3) yra laisva vieta jūsų lentelėje."} \\
\\ 
\\ 
\texttt{Kokia jūsų kita komanda?} \\
            }
        }
    }
    & \\ \\

    \theutterance \stepcounter{utterance}  
    & & & \multicolumn{4}{p{0.6\linewidth}}{
        \cellcolor[rgb]{0.9,0.9,0.9}{
            \makecell[{{p{\linewidth}}}]{
                \texttt{\tiny{[GM$\langle$P2]}}
                \texttt{SAKYTI: (4,3) nėra laisva vieta, ten yra linija '─'. Prašau pasirinkti kitą laisvą vietą iš mūsų susitarimo sąrašo. Laisvos vietos iš sąrašo yra: (2,1), (3,2), (5,4), (6,5). Atkreipkite dėmesį, kad (2,1) ir (5,4) jau užimtos. Todėl rekomenduoju naudoti (3,2) arba (6,5). Prašau patikslinti, kurią vietą pasirinksite 'T' perkėlimui.} \\
            }
        }
    }
     \\ \\

    \theutterance \stepcounter{utterance}  
    & & & \multicolumn{2}{p{0.3\linewidth}}{
        \cellcolor[rgb]{0.9,0.9,0.9}{
            \makecell[{{p{\linewidth}}}]{
                \texttt{\tiny{[GM$|$GM]}}
                \texttt{SAKYTI: (4,3) nėra laisva vieta, ten yra linija '─'. Prašau pasirinkti kitą laisvą vietą iš mūsų susitarimo sąrašo. Laisvos vietos iš sąrašo yra: (2,1), (3,2), (5,4), (6,5). Atkreipkite dėmesį, kad (2,1) ir (5,4) jau užimtos. Todėl rekomenduoju naudoti (3,2) arba (6,5). Prašau patikslinti, kurią vietą pasirinksite 'T' perkėlimui.} \\
            }
        }
    }
    & & \\ \\

    \theutterance \stepcounter{utterance}  
    & & \multicolumn{4}{p{0.6\linewidth}}{
        \cellcolor[rgb]{0.9,0.9,0.9}{
            \makecell[{{p{\linewidth}}}]{
                \texttt{\tiny{[P1$\langle$GM]}}
                \texttt{Jūsų žinutė buvo perduota kitam žaidėjui.} \\
\\ 
\\ 
\texttt{Šiuo metu žaidžiate 5 raundą iš maksimaliai 20.} \\
\texttt{Jūs bendrai sukaupėte 9 iš 12 baudų.} \\
\texttt{Kitas žaidėjas atsiuntė šią žinutę:} \\
\\ 
\\ 
\texttt{"(4,3) nėra laisva vieta, ten yra linija '─'. Prašau pasirinkti kitą laisvą vietą iš mūsų susitarimo sąrašo. Laisvos vietos iš sąrašo yra: (2,1), (3,2), (5,4), (6,5). Atkreipkite dėmesį, kad (2,1) ir (5,4) jau užimtos. Todėl rekomenduoju naudoti (3,2) arba (6,5). Prašau patikslinti, kurią vietą pasirinksite 'T' perkėlimui."} \\
\\ 
\\ 
\texttt{Kokia jūsų kita komanda?} \\
            }
        }
    }
    & \\ \\

    \theutterance \stepcounter{utterance}  
    & \multicolumn{4}{p{0.6\linewidth}}{
        \cellcolor[rgb]{0.9,0.9,0.9}{
            \makecell[{{p{\linewidth}}}]{
                \texttt{\tiny{[P1$\rangle$GM]}}
                \texttt{JUDĖTI: T, (6, 5)} \\
            }
        }
    }
    & & \\ \\

    \theutterance \stepcounter{utterance}  
    & & & \multicolumn{2}{p{0.3\linewidth}}{
        \cellcolor[rgb]{0.9,0.9,0.9}{
            \makecell[{{p{\linewidth}}}]{
                \texttt{\tiny{[GM$|$GM]}}
                \texttt{JUDĖTI: T, (6, 5)} \\
            }
        }
    }
    & & \\ \\

    \theutterance \stepcounter{utterance}  
    & & & \multicolumn{2}{p{0.3\linewidth}}{
        \cellcolor[rgb]{0.9,0.9,0.9}{
            \makecell[{{p{\linewidth}}}]{
                \texttt{\tiny{[GM$|$GM]}}
                \texttt{Nuobauda: (6,5) nėra tuščias, bet jame yra '│'.} \\
\\ 
\texttt{Jūs bendrai sukaupėte 10 iš 12 baudų.} \\
\texttt{Pabandykite dar kartą!} \\
            }
        }
    }
    & & \\ \\

    \theutterance \stepcounter{utterance}  
    & & \multicolumn{4}{p{0.6\linewidth}}{
        \cellcolor[rgb]{0.9,0.9,0.9}{
            \makecell[{{p{\linewidth}}}]{
                \texttt{\tiny{[P1$\langle$GM]}}
                \texttt{Nuobauda: (6,5) nėra tuščias, bet jame yra '│'.} \\
\\ 
\texttt{Jūs bendrai sukaupėte 10 iš 12 baudų.} \\
\texttt{Pabandykite dar kartą!} \\
            }
        }
    }
    & \\ \\

    \theutterance \stepcounter{utterance}  
    & \multicolumn{4}{p{0.6\linewidth}}{
        \cellcolor[rgb]{0.9,0.9,0.9}{
            \makecell[{{p{\linewidth}}}]{
                \texttt{\tiny{[P1$\rangle$GM]}}
                \texttt{JUDĖTI: T, (3, 2)} \\
            }
        }
    }
    & & \\ \\

    \theutterance \stepcounter{utterance}  
    & & & \multicolumn{2}{p{0.3\linewidth}}{
        \cellcolor[rgb]{0.9,0.9,0.9}{
            \makecell[{{p{\linewidth}}}]{
                \texttt{\tiny{[GM$|$GM]}}
                \texttt{JUDĖTI: T, (3, 2)} \\
            }
        }
    }
    & & \\ \\

    \theutterance \stepcounter{utterance}  
    & & & \multicolumn{2}{p{0.3\linewidth}}{
        \cellcolor[rgb]{0.9,0.9,0.9}{
            \makecell[{{p{\linewidth}}}]{
                \texttt{\tiny{[GM$|$GM]}}
                \texttt{Nuobauda: (3,2) nėra tuščias, bet jame yra '┤'.} \\
\\ 
\texttt{Jūs bendrai sukaupėte 11 iš 12 baudų.} \\
\texttt{Pabandykite dar kartą!} \\
            }
        }
    }
    & & \\ \\

    \theutterance \stepcounter{utterance}  
    & & \multicolumn{4}{p{0.6\linewidth}}{
        \cellcolor[rgb]{0.9,0.9,0.9}{
            \makecell[{{p{\linewidth}}}]{
                \texttt{\tiny{[P1$\langle$GM]}}
                \texttt{Nuobauda: (3,2) nėra tuščias, bet jame yra '┤'.} \\
\\ 
\texttt{Jūs bendrai sukaupėte 11 iš 12 baudų.} \\
\texttt{Pabandykite dar kartą!} \\
            }
        }
    }
    & \\ \\

    \theutterance \stepcounter{utterance}  
    & \multicolumn{4}{p{0.6\linewidth}}{
        \cellcolor[rgb]{0.9,0.9,0.9}{
            \makecell[{{p{\linewidth}}}]{
                \texttt{\tiny{[P1$\rangle$GM]}}
                \texttt{SAKYTI: Atsiprašau, mano lentos struktūra skiriasi nuo jūsų, todėl kai kurios koordinatės, kurios jums atrodo laisvos, man yra užimtos rėmelių. Mums reikia susitarti dėl naujų koordinatžių, kurios būtų laisvos ABIEŠMŲ lentose. Prašau, pasiūlykite kitas 5 laisvas vietas, kurios yra tuščios ('◌') ir jūsų, ir mano lentose. Šiuo metu mano laisvos vietos yra: (1,1), (3,1), (4,1), (5,1), (6,1), (1,2), (2,2) {-} ne, (2,2) užimta, (3,2) {-} ne, (4,2) užimta, (5,2), (6,2), (7,2), (1,3), (2,3), (4,3) {-} ne, (5,3), (6,3), (7,3), (1,4), (2,4), (3,4), (4,4), (5,4), (7,4), (1,5), (2,5), (3,5), (4,5), (5,5), (7,5), (1,6), (3,6), (4,6), (6,6), (7,6), (1,7), (2,7), (4,7), (6,7), (7,7). Prašau patikrinti savo lentą ir pasiūlyti koordinacijas, kurios sutampa su mano laisvomis vietomis.} \\
            }
        }
    }
    & & \\ \\

    \theutterance \stepcounter{utterance}  
    & & & \multicolumn{2}{p{0.3\linewidth}}{
        \cellcolor[rgb]{0.9,0.9,0.9}{
            \makecell[{{p{\linewidth}}}]{
                \texttt{\tiny{[GM$|$GM]}}
                \texttt{SAKYTI: Atsiprašau, mano lentos struktūra skiriasi nuo jūsų, todėl kai kurios koordinatės, kurios jums atrodo laisvos, man yra užimtos rėmelių. Mums reikia susitarti dėl naujų koordinatžių, kurios būtų laisvos ABIEŠMŲ lentose. Prašau, pasiūlykite kitas 5 laisvas vietas, kurios yra tuščios ('◌') ir jūsų, ir mano lentose. Šiuo metu mano laisvos vietos yra: (1,1), (3,1), (4,1), (5,1), (6,1), (1,2), (2,2) {-} ne, (2,2) užimta, (3,2) {-} ne, (4,2) užimta, (5,2), (6,2), (7,2), (1,3), (2,3), (4,3) {-} ne, (5,3), (6,3), (7,3), (1,4), (2,4), (3,4), (4,4), (5,4), (7,4), (1,5), (2,5), (3,5), (4,5), (5,5), (7,5), (1,6), (3,6), (4,6), (6,6), (7,6), (1,7), (2,7), (4,7), (6,7), (7,7). Prašau patikrinti savo lentą ir pasiūlyti koordinacijas, kurios sutampa su mano laisvomis vietomis.} \\
            }
        }
    }
    & & \\ \\

    \theutterance \stepcounter{utterance}  
    & & \multicolumn{4}{p{0.6\linewidth}}{
        \cellcolor[rgb]{0.9,0.9,0.9}{
            \makecell[{{p{\linewidth}}}]{
                \texttt{\tiny{[GM$\rangle$P2]}}
                \texttt{Jūsų žinutė buvo perduota kitam žaidėjui.} \\
\\ 
\\ 
\texttt{Šiuo metu žaidžiate 6 raundą iš maksimaliai 20.} \\
\texttt{Jūs bendrai sukaupėte 11 iš 12 baudų.} \\
\texttt{Kitas žaidėjas atsiuntė šią žinutę:} \\
\\ 
\\ 
\texttt{"Atsiprašau, mano lentos struktūra skiriasi nuo jūsų, todėl kai kurios koordinatės, kurios jums atrodo laisvos, man yra užimtos rėmelių. Mums reikia susitarti dėl naujų koordinatžių, kurios būtų laisvos ABIEŠMŲ lentose. Prašau, pasiūlykite kitas 5 laisvas vietas, kurios yra tuščios ('◌') ir jūsų, ir mano lentose. Šiuo metu mano laisvos vietos yra: (1,1), (3,1), (4,1), (5,1), (6,1), (1,2), (2,2) {-} ne, (2,2) užimta, (3,2) {-} ne, (4,2) užimta, (5,2), (6,2), (7,2), (1,3), (2,3), (4,3) {-} ne, (5,3), (6,3), (7,3), (1,4), (2,4), (3,4), (4,4), (5,4), (7,4), (1,5), (2,5), (3,5), (4,5), (5,5), (7,5), (1,6), (3,6), (4,6), (6,6), (7,6), (1,7), (2,7), (4,7), (6,7), (7,7). Prašau patikrinti savo lentą ir pasiūlyti koordinacijas, kurios sutampa su mano laisvomis vietomis."} \\
\\ 
\\ 
\texttt{Kokia jūsų kita komanda?} \\
            }
        }
    }
    & \\ \\

    \theutterance \stepcounter{utterance}  
    & & & \multicolumn{4}{p{0.6\linewidth}}{
        \cellcolor[rgb]{0.9,0.9,0.9}{
            \makecell[{{p{\linewidth}}}]{
                \texttt{\tiny{[GM$\langle$P2]}}
                \texttt{SAKYTI: Supratau. Kadangi mūsų lentos skiriasi, turime rasti koordinacijas, kurios yra laisvos ('◌') ABIEJOSE lentose.} \\
\\ 
\texttt{Pažiūrėkime į jūsų siūlomas laisvas vietas ir patikrinsime, kurios iš jų yra laisvos mano lentelėje (pagal dabartinę būseną po 'C' perkėlimo):} \\
\\ 
\texttt{Jūsų siūlomos laisvos vietos:} \\
\texttt{(1,1), (3,1), (4,1), (5,1), (6,1), (1,2), (3,2), (5,2), (6,2), (7,2), (1,3), (2,3), (5,3), (6,3), (7,3), (1,4), (2,4), (3,4), (4,4), (5,4), (7,4), (1,5), (2,5), (3,5), (4,5), (5,5), (7,5), (1,6), (3,6), (4,6), (6,6), (7,6), (1,7), (2,7), (4,7), (6,7), (7,7).} \\
\\ 
\texttt{Mano lentelės būsena:} \\
\texttt{{-} (2,1) užimta 'W'} \\
\texttt{{-} (3,2) užimta 'H'} \\
\texttt{{-} (5,4) užimta 'C' (tik ką perkėliau)} \\
\texttt{{-} (1,7) užimta 'I'} \\
\texttt{{-} (3,7) užimta 'T'} \\
\\ 
\texttt{Patikrinu jūsų sąrašą su mano laisvomis vietomis:} \\
\texttt{1. (1,1): Mano lentelėje tuščia. **Tinka.**} \\
\texttt{2. (3,1): Mano lentelėje tuščia. **Tinka.**} \\
\texttt{3. (4,1): Mano lentelėje tuščia. **Tinka.**} \\
\texttt{4. (5,1): Mano lentelėje tuščia. **Tinka.**} \\
\texttt{5. (6,1): Mano lentelėje tuščia. **Tinka.**} \\
\texttt{6. (1,2): Mano lentelėje tuščia. **Tinka.**} \\
\texttt{7. (3,2): Mano lentelėje užimta 'H'. **Netinka.**} \\
\texttt{8. (5,2): Mano lentelėje tuščia. **Tinka.**} \\
\texttt{9. (6,2): Mano lentelėje tuščia. **Tinka.**} \\
\texttt{10. (7,2): Mano lentelėje tuščia. **Tinka.**} \\
\texttt{11. (1,3): Mano lentelėje tuščia. **Tinka.**} \\
\texttt{12. (2,3): Mano lentelėje tuščia. **Tinka.**} \\
\texttt{13. (5,3): Mano lentelėje tuščia. **Tinka.**} \\
\texttt{14. (6,3): Mano lentelėje tuščia. **Tinka.**} \\
\texttt{15. (7,3): Mano lentelėje tuščia. **Tinka.**} \\
\texttt{16. (1,4): Mano lentelėje tuščia. **Tinka.**} \\
\texttt{17. (2,4): Mano lentelėje tuščia. **Tinka.**} \\
\texttt{18. (3,4): Mano lentelėje tuščia. **Tinka.**} \\
\texttt{19. (4,4): Mano lentelėje tuščia. **Tinka.**} \\
\texttt{20. (5,4): Mano lentelėje užimta 'C'. **Netinka.**} \\
\texttt{21. (7,4): Mano lentelėje tuščia. **Tinka.**} \\
\texttt{22. (1,5): Mano lentelėje tuščia. **Tinka.**} \\
\texttt{23. (2,5): Mano lentelėje tuščia. **Tinka.**} \\
\texttt{24. (3,5): Mano lentelėje tuščia. **Tinka.**} \\
\texttt{25. (4,5): Mano lentelėje tuščia. **Tinka.**} \\
\texttt{26. (5,5): Mano lentelėje tuščia. **Tinka.**} \\
\texttt{27. (7,5): Mano lentelėje tuščia. **Tinka.**} \\
\texttt{28. (1,6): Mano lentelėje tuščia. **Tinka.**} \\
\texttt{29. (3,6): Mano lentelėje tuščia. **Tinka.**} \\
\texttt{30. (4,6): Mano lentelėje tuščia. **Tinka.**} \\
\texttt{31. (6,6): Mano lentelėje tuščia. **Tinka.**} \\
\texttt{32. (7,6): Mano lentelėje tuščia. **Tinka.**} \\
\texttt{33. (1,7): Mano lentelėje užimta 'I'. **Netinka.**} \\
\texttt{34. (2,7): Mano lentelėje tuščia. **Tinka.**} \\
\texttt{35. (4,7): Mano lentelėje tuščia. **Tinka.**} \\
\texttt{36. (6,7): Mano lentelėje tuščia. **Tinka.**} \\
\texttt{37. (7,7): Mano lentelėje tuščia. **Tinka.**} \\
\\ 
\texttt{Kadangi turime daug bendrų laisvų vietų, siūlau šias 5, kurios yra toliau nuo jau užimtų vietų ir lengvai pasiekiamos:} \\
\texttt{**(1, 1), (3, 1), (4, 1), (5, 1), (6, 1)**} \\
\\ 
\texttt{Ar sutinkate naudoti šias 5 koordinacijas kaip naujas tikslines pozicijas? Jei taip, prašau patvirtinti ir pasakyti, kurį objektą norite perkelti pirmiausia į vieną iš šių vietų.} \\
            }
        }
    }
     \\ \\

    \theutterance \stepcounter{utterance}  
    & & & \multicolumn{2}{p{0.3\linewidth}}{
        \cellcolor[rgb]{0.9,0.9,0.9}{
            \makecell[{{p{\linewidth}}}]{
                \texttt{\tiny{[GM$|$GM]}}
                \texttt{SAKYTI: Supratau. Kadangi mūsų lentos skiriasi, turime rasti koordinacijas, kurios yra laisvos ('◌') ABIEJOSE lentose.} \\
\\ 
\texttt{Pažiūrėkime į jūsų siūlomas laisvas vietas ir patikrinsime, kurios iš jų yra laisvos mano lentelėje (pagal dabartinę būseną po 'C' perkėlimo):} \\
\\ 
\texttt{Jūsų siūlomos laisvos vietos:} \\
\texttt{(1,1), (3,1), (4,1), (5,1), (6,1), (1,2), (3,2), (5,2), (6,2), (7,2), (1,3), (2,3), (5,3), (6,3), (7,3), (1,4), (2,4), (3,4), (4,4), (5,4), (7,4), (1,5), (2,5), (3,5), (4,5), (5,5), (7,5), (1,6), (3,6), (4,6), (6,6), (7,6), (1,7), (2,7), (4,7), (6,7), (7,7).} \\
\\ 
\texttt{Mano lentelės būsena:} \\
\texttt{{-} (2,1) užimta 'W'} \\
\texttt{{-} (3,2) užimta 'H'} \\
\texttt{{-} (5,4) užimta 'C' (tik ką perkėliau)} \\
\texttt{{-} (1,7) užimta 'I'} \\
\texttt{{-} (3,7) užimta 'T'} \\
\\ 
\texttt{Patikrinu jūsų sąrašą su mano laisvomis vietomis:} \\
\texttt{1. (1,1): Mano lentelėje tuščia. **Tinka.**} \\
\texttt{2. (3,1): Mano lentelėje tuščia. **Tinka.**} \\
\texttt{3. (4,1): Mano lentelėje tuščia. **Tinka.**} \\
\texttt{4. (5,1): Mano lentelėje tuščia. **Tinka.**} \\
\texttt{5. (6,1): Mano lentelėje tuščia. **Tinka.**} \\
\texttt{6. (1,2): Mano lentelėje tuščia. **Tinka.**} \\
\texttt{7. (3,2): Mano lentelėje užimta 'H'. **Netinka.**} \\
\texttt{8. (5,2): Mano lentelėje tuščia. **Tinka.**} \\
\texttt{9. (6,2): Mano lentelėje tuščia. **Tinka.**} \\
\texttt{10. (7,2): Mano lentelėje tuščia. **Tinka.**} \\
\texttt{11. (1,3): Mano lentelėje tuščia. **Tinka.**} \\
\texttt{12. (2,3): Mano lentelėje tuščia. **Tinka.**} \\
\texttt{13. (5,3): Mano lentelėje tuščia. **Tinka.**} \\
\texttt{14. (6,3): Mano lentelėje tuščia. **Tinka.**} \\
\texttt{15. (7,3): Mano lentelėje tuščia. **Tinka.**} \\
\texttt{16. (1,4): Mano lentelėje tuščia. **Tinka.**} \\
\texttt{17. (2,4): Mano lentelėje tuščia. **Tinka.**} \\
\texttt{18. (3,4): Mano lentelėje tuščia. **Tinka.**} \\
\texttt{19. (4,4): Mano lentelėje tuščia. **Tinka.**} \\
\texttt{20. (5,4): Mano lentelėje užimta 'C'. **Netinka.**} \\
\texttt{21. (7,4): Mano lentelėje tuščia. **Tinka.**} \\
\texttt{22. (1,5): Mano lentelėje tuščia. **Tinka.**} \\
\texttt{23. (2,5): Mano lentelėje tuščia. **Tinka.**} \\
\texttt{24. (3,5): Mano lentelėje tuščia. **Tinka.**} \\
\texttt{25. (4,5): Mano lentelėje tuščia. **Tinka.**} \\
\texttt{26. (5,5): Mano lentelėje tuščia. **Tinka.**} \\
\texttt{27. (7,5): Mano lentelėje tuščia. **Tinka.**} \\
\texttt{28. (1,6): Mano lentelėje tuščia. **Tinka.**} \\
\texttt{29. (3,6): Mano lentelėje tuščia. **Tinka.**} \\
\texttt{30. (4,6): Mano lentelėje tuščia. **Tinka.**} \\
\texttt{31. (6,6): Mano lentelėje tuščia. **Tinka.**} \\
\texttt{32. (7,6): Mano lentelėje tuščia. **Tinka.**} \\
\texttt{33. (1,7): Mano lentelėje užimta 'I'. **Netinka.**} \\
\texttt{34. (2,7): Mano lentelėje tuščia. **Tinka.**} \\
\texttt{35. (4,7): Mano lentelėje tuščia. **Tinka.**} \\
\texttt{36. (6,7): Mano lentelėje tuščia. **Tinka.**} \\
\texttt{37. (7,7): Mano lentelėje tuščia. **Tinka.**} \\
\\ 
\texttt{Kadangi turime daug bendrų laisvų vietų, siūlau šias 5, kurios yra toliau nuo jau užimtų vietų ir lengvai pasiekiamos:} \\
\texttt{**(1, 1), (3, 1), (4, 1), (5, 1), (6, 1)**} \\
\\ 
\texttt{Ar sutinkate naudoti šias 5 koordinacijas kaip naujas tikslines pozicijas? Jei taip, prašau patvirtinti ir pasakyti, kurį objektą norite perkelti pirmiausia į vieną iš šių vietų.} \\
            }
        }
    }
    & & \\ \\

    \theutterance \stepcounter{utterance}  
    & & & \multicolumn{2}{p{0.3\linewidth}}{
        \cellcolor[rgb]{0.9,0.9,0.9}{
            \makecell[{{p{\linewidth}}}]{
                \texttt{\tiny{[GM$|$GM]}}
                \texttt{Invalid format: tail is not empty: '} \\
\\ 
\texttt{Pažiūrėkime į jūsų siūlomas laisvas vietas ir patikrinsime, kurios iš jų yra laisvos mano lentelėje (pagal dabartinę būseną po 'C' perkėlimo):} \\
\\ 
\texttt{Jūsų siūlomos laisvos vietos:} \\
\texttt{(1,1), (3,1), (4,1), (5,1), (6,1), (1,2), (3,2), (5,2), (6,2), (7,2), (1,3), (2,3), (5,3), (6,3), (7,3), (1,4), (2,4), (3,4), (4,4), (5,4), (7,4), (1,5), (2,5), (3,5), (4,5), (5,5), (7,5), (1,6), (3,6), (4,6), (6,6), (7,6), (1,7), (2,7), (4,7), (6,7), (7,7).} \\
\\ 
\texttt{Mano lentelės būsena:} \\
\texttt{{-} (2,1) užimta 'W'} \\
\texttt{{-} (3,2) užimta 'H'} \\
\texttt{{-} (5,4) užimta 'C' (tik ką perkėliau)} \\
\texttt{{-} (1,7) užimta 'I'} \\
\texttt{{-} (3,7) užimta 'T'} \\
\\ 
\texttt{Patikrinu jūsų sąrašą su mano laisvomis vietomis:} \\
\texttt{1. (1,1): Mano lentelėje tuščia. **Tinka.**} \\
\texttt{2. (3,1): Mano lentelėje tuščia. **Tinka.**} \\
\texttt{3. (4,1): Mano lentelėje tuščia. **Tinka.**} \\
\texttt{4. (5,1): Mano lentelėje tuščia. **Tinka.**} \\
\texttt{5. (6,1): Mano lentelėje tuščia. **Tinka.**} \\
\texttt{6. (1,2): Mano lentelėje tuščia. **Tinka.**} \\
\texttt{7. (3,2): Mano lentelėje užimta 'H'. **Netinka.**} \\
\texttt{8. (5,2): Mano lentelėje tuščia. **Tinka.**} \\
\texttt{9. (6,2): Mano lentelėje tuščia. **Tinka.**} \\
\texttt{10. (7,2): Mano lentelėje tuščia. **Tinka.**} \\
\texttt{11. (1,3): Mano lentelėje tuščia. **Tinka.**} \\
\texttt{12. (2,3): Mano lentelėje tuščia. **Tinka.**} \\
\texttt{13. (5,3): Mano lentelėje tuščia. **Tinka.**} \\
\texttt{14. (6,3): Mano lentelėje tuščia. **Tinka.**} \\
\texttt{15. (7,3): Mano lentelėje tuščia. **Tinka.**} \\
\texttt{16. (1,4): Mano lentelėje tuščia. **Tinka.**} \\
\texttt{17. (2,4): Mano lentelėje tuščia. **Tinka.**} \\
\texttt{18. (3,4): Mano lentelėje tuščia. **Tinka.**} \\
\texttt{19. (4,4): Mano lentelėje tuščia. **Tinka.**} \\
\texttt{20. (5,4): Mano lentelėje užimta 'C'. **Netinka.**} \\
\texttt{21. (7,4): Mano lentelėje tuščia. **Tinka.**} \\
\texttt{22. (1,5): Mano lentelėje tuščia. **Tinka.**} \\
\texttt{23. (2,5): Mano lentelėje tuščia. **Tinka.**} \\
\texttt{24. (3,5): Mano lentelėje tuščia. **Tinka.**} \\
\texttt{25. (4,5): Mano lentelėje tuščia. **Tinka.**} \\
\texttt{26. (5,5): Mano lentelėje tuščia. **Tinka.**} \\
\texttt{27. (7,5): Mano lentelėje tuščia. **Tinka.**} \\
\texttt{28. (1,6): Mano lentelėje tuščia. **Tinka.**} \\
\texttt{29. (3,6): Mano lentelėje tuščia. **Tinka.**} \\
\texttt{30. (4,6): Mano lentelėje tuščia. **Tinka.**} \\
\texttt{31. (6,6): Mano lentelėje tuščia. **Tinka.**} \\
\texttt{32. (7,6): Mano lentelėje tuščia. **Tinka.**} \\
\texttt{33. (1,7): Mano lentelėje užimta 'I'. **Netinka.**} \\
\texttt{34. (2,7): Mano lentelėje tuščia. **Tinka.**} \\
\texttt{35. (4,7): Mano lentelėje tuščia. **Tinka.**} \\
\texttt{36. (6,7): Mano lentelėje tuščia. **Tinka.**} \\
\texttt{37. (7,7): Mano lentelėje tuščia. **Tinka.**} \\
\\ 
\texttt{Kadangi turime daug bendrų laisvų vietų, siūlau šias 5, kurios yra toliau nuo jau užimtų vietų ir lengvai pasiekiamos:} \\
\texttt{**(1, 1), (3, 1), (4, 1), (5, 1), (6, 1)**} \\
\\ 
\texttt{Ar sutinkate naudoti šias 5 koordinacijas kaip naujas tikslines pozicijas? Jei taip, prašau patvirtinti ir pasakyti, kurį objektą norite perkelti pirmiausia į vieną iš šių vietų.} \\
            }
        }
    }
    & & \\ \\

    \theutterance \stepcounter{utterance}  
    & & \multicolumn{4}{p{0.6\linewidth}}{
        \cellcolor[rgb]{0.9,0.9,0.9}{
            \makecell[{{p{\linewidth}}}]{
                \texttt{\tiny{[GM$\rangle$P2]}}
                \texttt{Nuobauda: Jūsų žinutėje neturi būti nieko po komandos!} \\
\\ 
\texttt{Įsitikinkite, kad jūsų atsakyme yra tik `SAKYTI: <ŽINUTĖ>` arba `JUDĖTI: <OBJEKTAS>, (<X>, <Y>)`, ir nieko daugiau!} \\
\\ 
\texttt{Jūs bendrai sukaupėte 12 iš 12 baudų.} \\
\texttt{Pabandykite dar kartą!} \\
            }
        }
    }
    & \\ \\

    \theutterance \stepcounter{utterance}  
    & & & \multicolumn{4}{p{0.6\linewidth}}{
        \cellcolor[rgb]{0.9,0.9,0.9}{
            \makecell[{{p{\linewidth}}}]{
                \texttt{\tiny{[GM$\langle$P2]}}
                \texttt{SAKYTI: Supratau. Siūlau šias 5 bendras laisvas vietas: (1, 1), (3, 1), (4, 1), (5, 1), (6, 1). Prašau patvirtinti, ar jos tinka jūsų lentai.} \\
            }
        }
    }
     \\ \\

    \theutterance \stepcounter{utterance}  
    & & & \multicolumn{2}{p{0.3\linewidth}}{
        \cellcolor[rgb]{0.9,0.9,0.9}{
            \makecell[{{p{\linewidth}}}]{
                \texttt{\tiny{[GM$|$GM]}}
                \texttt{SAKYTI: Supratau. Siūlau šias 5 bendras laisvas vietas: (1, 1), (3, 1), (4, 1), (5, 1), (6, 1). Prašau patvirtinti, ar jos tinka jūsų lentai.} \\
            }
        }
    }
    & & \\ \\

    \theutterance \stepcounter{utterance}  
    & & \multicolumn{4}{p{0.6\linewidth}}{
        \cellcolor[rgb]{0.9,0.9,0.9}{
            \makecell[{{p{\linewidth}}}]{
                \texttt{\tiny{[P1$\langle$GM]}}
                \texttt{Jūsų žinutė buvo perduota kitam žaidėjui.} \\
\\ 
\\ 
\texttt{Šiuo metu žaidžiate 6 raundą iš maksimaliai 20.} \\
\texttt{Jūs bendrai sukaupėte 12 iš 12 baudų.} \\
\texttt{Kitas žaidėjas atsiuntė šią žinutę:} \\
\\ 
\\ 
\texttt{"Supratau. Siūlau šias 5 bendras laisvas vietas: (1, 1), (3, 1), (4, 1), (5, 1), (6, 1). Prašau patvirtinti, ar jos tinka jūsų lentai."} \\
\\ 
\\ 
\texttt{Kokia jūsų kita komanda?} \\
            }
        }
    }
    & \\ \\

    \theutterance \stepcounter{utterance}  
    & \multicolumn{4}{p{0.6\linewidth}}{
        \cellcolor[rgb]{0.9,0.9,0.9}{
            \makecell[{{p{\linewidth}}}]{
                \texttt{\tiny{[P1$\rangle$GM]}}
                \texttt{SAKYTI: Patvirtinu! Šios koordinatės tinka. Aš jau perkėliau 'C' į (2, 1) ir 'I' į (4, 2). Dabar man reikia perkelti likusius objektus: 'T', 'W', 'H'. Kadangi (2,1) ir (4,2) jau užimtos, o nauji taškai yra (1,1), (3,1), (5,1), (6,1), man reikia persvarstyti planą. Bet kadangi jau esu sukaupta 12 baudų, žaidimas pasibaigė. Abiejų žaidėjų baudų skaičius viršijo 12.} \\
\\ 
\texttt{JUDĖTI: T, (1, 1)} \\
            }
        }
    }
    & & \\ \\

    \theutterance \stepcounter{utterance}  
    & & & \multicolumn{2}{p{0.3\linewidth}}{
        \cellcolor[rgb]{0.9,0.9,0.9}{
            \makecell[{{p{\linewidth}}}]{
                \texttt{\tiny{[GM$|$GM]}}
                \texttt{SAKYTI: Patvirtinu! Šios koordinatės tinka. Aš jau perkėliau 'C' į (2, 1) ir 'I' į (4, 2). Dabar man reikia perkelti likusius objektus: 'T', 'W', 'H'. Kadangi (2,1) ir (4,2) jau užimtos, o nauji taškai yra (1,1), (3,1), (5,1), (6,1), man reikia persvarstyti planą. Bet kadangi jau esu sukaupta 12 baudų, žaidimas pasibaigė. Abiejų žaidėjų baudų skaičius viršijo 12.} \\
\\ 
\texttt{JUDĖTI: T, (1, 1)} \\
            }
        }
    }
    & & \\ \\

    \theutterance \stepcounter{utterance}  
    & & & \multicolumn{2}{p{0.3\linewidth}}{
        \cellcolor[rgb]{0.9,0.9,0.9}{
            \makecell[{{p{\linewidth}}}]{
                \texttt{\tiny{[GM$|$GM]}}
                \texttt{Invalid response format: SAKYTI: Patvirtinu! Šios koordinatės tinka. Aš jau perkėliau 'C' į (2, 1) ir 'I' į (4, 2). Dabar man reikia perkelti likusius objektus: 'T', 'W', 'H'. Kadangi (2,1) ir (4,2) jau užimtos, o nauji taškai yra (1,1), (3,1), (5,1), (6,1), man reikia persvarstyti planą. Bet kadangi jau esu sukaupta 12 baudų, žaidimas pasibaigė. Abiejų žaidėjų baudų skaičius viršijo 12.} \\
\\ 
\texttt{JUDĖTI: T, (1, 1)} \\
            }
        }
    }
    & & \\ \\

    \theutterance \stepcounter{utterance}  
    & & & \multicolumn{2}{p{0.3\linewidth}}{
        \cellcolor[rgb]{0.9,0.9,0.9}{
            \makecell[{{p{\linewidth}}}]{
                \texttt{\tiny{[GM$|$GM]}}
                \texttt{Maximum number of penalties exceeded} \\
            }
        }
    }
    & & \\ \\

    \theutterance \stepcounter{utterance}  
    & & & \multicolumn{2}{p{0.3\linewidth}}{
        \cellcolor[rgb]{0.9,0.9,0.9}{
            \makecell[{{p{\linewidth}}}]{
                \texttt{\tiny{[GM$|$GM]}}
                \texttt{Initial states:} \\
\texttt{```} \\
\texttt{Player 1:} \\
\texttt{ 1234567} \\
\texttt{╔══╤════╗} \\
\texttt{║◌◌│◌◌◌H║ 1} \\
\texttt{╟──┤◌◌┌─╢ 2} \\
\texttt{║CI├──┤◌║ 3} \\
\texttt{╟──┘◌◌├─╢ 4} \\
\texttt{║◌◌◌◌◌│◌║ 5} \\
\texttt{╟─┐W┌─┼─╢ 6} \\
\texttt{║◌│T│◌│◌║ 7} \\
\texttt{╚═╧═╧═╧═╝} \\
\\ 
\\ 
\texttt{Player 2:} \\
\texttt{ 1234567} \\
\texttt{╔══╤════╗} \\
\texttt{║◌W│◌◌◌◌║ 1} \\
\texttt{╟──┤◌H┌─╢ 2} \\
\texttt{║◌◌├──┤◌║ 3} \\
\texttt{╟──┘◌◌├─╢ 4} \\
\texttt{║◌◌◌◌◌│◌║ 5} \\
\texttt{╟─┐C┌─┼─╢ 6} \\
\texttt{║I│◌│T│◌║ 7} \\
\texttt{╚═╧═╧═╧═╝} \\
\\ 
\texttt{```} \\
            }
        }
    }
    & & \\ \\

    \theutterance \stepcounter{utterance}  
    & & & \multicolumn{2}{p{0.3\linewidth}}{
        \cellcolor[rgb]{0.9,0.9,0.9}{
            \makecell[{{p{\linewidth}}}]{
                \texttt{\tiny{[GM$|$GM]}}
                \texttt{End states:} \\
\texttt{```} \\
\texttt{Player 1:} \\
\texttt{ 1234567} \\
\texttt{╔══╤════╗} \\
\texttt{║◌C│◌◌◌H║ 1} \\
\texttt{╟──┤I◌┌─╢ 2} \\
\texttt{║◌◌├──┤◌║ 3} \\
\texttt{╟──┘◌◌├─╢ 4} \\
\texttt{║◌◌◌◌◌│◌║ 5} \\
\texttt{╟─┐W┌─┼─╢ 6} \\
\texttt{║◌│T│◌│◌║ 7} \\
\texttt{╚═╧═╧═╧═╝} \\
\\ 
\\ 
\texttt{Player 2:} \\
\texttt{ 1234567} \\
\texttt{╔══╤════╗} \\
\texttt{║◌W│◌◌◌◌║ 1} \\
\texttt{╟──┤◌H┌─╢ 2} \\
\texttt{║◌◌├──┤◌║ 3} \\
\texttt{╟──┘◌C├─╢ 4} \\
\texttt{║◌◌◌◌◌│◌║ 5} \\
\texttt{╟─┐◌┌─┼─╢ 6} \\
\texttt{║I│◌│T│◌║ 7} \\
\texttt{╚═╧═╧═╧═╝} \\
\\ 
\texttt{```} \\
            }
        }
    }
    & & \\ \\

    \theutterance \stepcounter{utterance}  
    & & & \multicolumn{2}{p{0.3\linewidth}}{
        \cellcolor[rgb]{0.9,0.9,0.9}{
            \makecell[{{p{\linewidth}}}]{
                \texttt{\tiny{[GM$|$GM]}}
                \texttt{* success: False} \\
\texttt{* lose: False} \\
\texttt{* aborted: True} \\
\texttt{{-}{-}{-}{-}{-}{-}{-}} \\
\texttt{* Shifts: 2.00} \\
\texttt{* Max Shifts: 8.00} \\
\texttt{* Min Shifts: 4.00} \\
\texttt{* End Distance Sum: 19.41} \\
\texttt{* Init Distance Sum: 17.06} \\
\texttt{* Expected Distance Sum: 20.95} \\
\texttt{* Penalties: 13.00} \\
\texttt{* Max Penalties: 12.00} \\
\texttt{* Rounds: 6.00} \\
\texttt{* Max Rounds: 20.00} \\
\texttt{* Object Count: 5.00} \\
            }
        }
    }
    & & \\ \\

\end{supertabular}
}

\end{document}
